% 3. 终端（命令行，命令提示符）
\section{终端（命令行，命令提示符）}

\begin{frame}
    \frametitle{终端是什么}
    终端有很多种，叫法也很多：
    \begin{itemize}
        \item Terminal
        \item Shell
        \item Command Line
        \item Command Prompt
        \item ……
    \end{itemize}
    本质上都是同一件东西：一个用文字命令和计算机对话的窗口。
\end{frame}

\begin{frame}
    \frametitle{打开终端}
    \begin{itemize}
        \item \alert{Linux（Ubuntu）}：按\code{Ctrl+Alt+T}，或在应用列表中搜索Terminal；
        \item \alert{Mac}：\code{Cmd+空格}搜索“终端”（Terminal）；
        \item \alert{Windows（WSL）}：开始菜单搜索Ubuntu，或打开PowerShell后输入\code{wsl}进入。
    \end{itemize}
\end{frame}

\begin{frame}
    \frametitle{在终端使用指令}
    bash常见命令（不一定全都掌握，但\code{cd}必须会）：
    \begin{itemize}
        \item \code{ls}：列出当前工作目录（Current Working Directory，CWD）下的文件；
        \begin{itemize}
            \item 默认情况下打开终端，工作目录是用户家目录\code{\textasciitilde}（Linux为\code{/home/用户名}，Mac为\code{/Users/用户名}）；
        \end{itemize}
        \item \code{cd somePath}：将工作目录切换到\code{somePath}，可以是相对或绝对路径；
        \begin{itemize}
            \item \code{cd ..}：切换到上一级目录；
            \item \code{cd 'My Folder'}：切换到\code{My Folder}，加引号是因为有空格；
        \end{itemize}
        \item \code{mkdir CS100}：在当前工作目录下创建名为\code{CS100}的文件夹；
        \item \code{rmdir CS100}：删除这个文件夹（要求文件夹为空）；
        \item \code{rm file}：删除名为\code{file}的文件；
        \item \code{rm -r CS100}：删除名为\code{CS100}的文件夹及其所有内容。
    \end{itemize}
\end{frame}

\begin{frame}
    \frametitle{验证编译器安装：PATH}
    输入\code{gcc}，shell会在\code{PATH}环境变量所包含的路径中寻找名为\code{gcc}的可执行文件，它通常位于\code{/usr/bin}。

    如果编译器的安装正确，应当看到如下输出：
    \begin{Code}*
        \lstinputlisting{code/terminal/gcc-no-input.txt}
    \end{Code}
    这不是报错，而是编译器在告诉你“给我点活干”。
\end{frame}

\begin{frame}
    \frametitle{查看编译器版本}
    给\code{gcc}加上\code{--version}参数即可查看版本：
    \begin{Code}*
        \lstinputlisting{code/terminal/gcc-version.sh}
    \end{Code}
    输出形如：
    \begin{Code}*
        \lstinputlisting{code/terminal/gcc-version-out.txt}
    \end{Code}

    顺带一提：\code{gcc}和\code{clang}是C编译器，\code{g++}和\code{clang++}是C++编译器。
\end{frame}

\begin{frame}
    \frametitle{编译运行一个Hello World程序}
    假设你的CS100课程的文件夹位于\code{\textasciitilde/courses/CS100}。

    在\code{\textasciitilde/courses/CS100/tmp}创建\code{hello.c}，用VSCode打开，输入如下代码：
    \begin{Code}*{c}
        \lstinputlisting{code/terminal/hello.c}
    \end{Code}
\end{frame}

\begin{frame}
    \frametitle{编译}
    打开终端并切换到\code{\textasciitilde/courses/CS100/tmp}，输入如下命令：
    \begin{Code}*
        \lstinputlisting{code/terminal/gcc-compile.sh}
    \end{Code}
    将会得到可执行文件\code{hello}。该文件名由紧跟在\code{-o}之后的字符串指定。
\end{frame}

\begin{frame}
    \frametitle{运行}
    在终端里输入\code{./hello}即可运行\code{hello}。

    \code{./}表示当前工作目录，因为Linux/Mac默认不会在当前目录下搜索可执行文件。

    输出：
    \begin{Code}*
        \lstinputlisting{code/terminal/hello-out.txt}
    \end{Code}
\end{frame}

\begin{frame}
    \frametitle{向编译器传递参数}
    \begin{Code}*
        \lstinputlisting{code/terminal/gcc-compile.sh}
    \end{Code}
    跟在\code{gcc}之后的各项是传递给\code{gcc}的\alert{参数}。

    \begin{itemize}
        \item \code{hello.c}是它需要编译的文件名，如果不在当前工作目录下，需要指明路径，例如\code{../../CS101/pa1/a.c}；
        \item \code{-o hello}指明生成的可执行文件的名字，也可以带有路径；Linux/Mac下可执行文件通常无后缀名；
        \item 还可以有其它参数，例如\code{-std=c17}指明采用2017年发布的语言标准C17，\code{-lm}链接数学库，\code{--save-temps}保留编译过程中产生的临时文件；
        \item 以上参数的顺序无所谓。
    \end{itemize}
\end{frame}
